% Part: propositional-logic
% Chapter: syntax-and-semantics
% Section: valuations-sat

\documentclass[../../../include/open-logic-section]{subfiles}

\begin{document}

\olfileid{pl}{syn}{val}

\olsection{\usetoken{P}{valuation} تے پورا اترنا}

\begin{defn}[!!^{valuation}s]
$\{\True, \False\}$ نوں دو صداقتی قدراں ``سچ'' تے ``جھوٹ'' دا
سیٹ منّو۔ $\Lang{L_0}$ لئی \emph{!!{valuation}} اک فنکشن~$\pAssign{v}$ اے
جیہڑا زبان دے !!{propositional variable}s نوں $\True$ یا $\False$
وچوں کوئی اک مقرر کردا اے، یعنی $\pAssign{v} \colon
\PVar \to \{\True, \False \}$۔
\end{defn}

\begin{defn}
  دتی ہوئی !!a{valuation}~$\pAssign{v}$ لئی، قدر معلوم کرن والے
  فنکشن $\pValue{v} \colon \Frm[L_0] \to \{\True, \False \}$ دی
  استقرائی تعریف ایویں کرو:
  \begin{align*}
    \iftag{prvFalse}{\pValue{v}(\lfalse) & = \False; \\ }{}
    \iftag{prvTrue}{\pValue{v}(\ltrue) & = \True; \\ }{}
    \pValue{v}(\Obj p_n) & = \pAssign{v}(\Obj p_n); \\
    \iftag{prvNot}{\pValue{v}(\lnot !A) & = \begin{cases}
        \True & \text{جے } \pValue{v}(!A) = \False;\\
        \False & \text{ہور صورت وچ۔}
      \end{cases} \\ }{}
    \iftag{prvAnd}{\pValue{v}(!A \land !B) & = \begin{cases}
        \True &
        \text{جے $\pValue{v}(!A) = \True$ تے $\pValue{v}(!B) = \True$؛}\\
        \False &
        \text{جے $\pValue{v}(!A) = \False$ یا $\pValue{v}(!B) = \False$}.
      \end{cases}\\ }{}
    \iftag{prvOr}{\pValue{v}(!A \lor !B) & = \begin{cases}
        \True &
        \text{جے $\pValue{v}(!A) = \True$ یا $\pValue{v}(!B) = \True$؛}\\
        \False &
        \text{جے $\pValue{v}(!A) = \False$ تے $\pValue{v}(!B) = \False$}.
      \end{cases}\\ }{}
    \iftag{prvIf}{\pValue{v}(!A \lif !B) & = \begin{cases}
        \True &
        \text{جے $\pValue{v}(!A) = \False$ یا $\pValue{v}(!B) = \True$؛}\\
        \False &
        \text{جے $\pValue{v}(!A) = \True$ تے $\pValue{v}(!B) = \False$}.
      \end{cases}\\ }{}
    \iftag{prvIff}{\pValue{v}(!A \liff !B) & = \begin{cases}
        \True &
        \text{جے $\pValue{v}(!A) = \pValue{v}(!B)$؛}\\
        \False &
        \text{جے $\pValue{v}(!A) \neq \pValue{v}(!B)$}.
      \end{cases}}{}
\end{align*}
\end{defn}

\begin{explain}
ایہ شرطاں اگلے صداقتی جدولاں دے مطابق نیں:
\begin{center}
\iftag{prvNot}{
  \begin{tabular}{|c||c|} \hline
    $!A$ & $ \lnot !A$ \\
    \hline \hline
    $\True$ & $\False$ \\
    $\False$ & $\True$ \\
    \hline
  \end{tabular}
}{}
\iftag{prvAnd}{
  \begin{tabular}{|cc||c|} \hline
    $!A$ & $!B$ & $!A \land !B$ \\
    \hline \hline
    $\True$ & $\True$ & $\True$ \\
    $\True$ & $\False$ & $\False$ \\
    $\False$ & $\True$ & $\False$ \\
    $\False$ & $\False$ & $\False$ \\
    \hline
  \end{tabular}
}{}
\iftag{prvOr}{
  \begin{tabular}{|cc||c|} \hline
    $!A$ & $!B$ & $!A \lor !B$ \\
    \hline \hline
    $\True$ & $\True$ & $\True$ \\
    $\True$ & $\False$ & $\True$ \\
    $\False$ & $\True$ & $\True$ \\
    $\False$ & $\False$ & $\False$ \\
    \hline
  \end{tabular}
}{}%
\iftag{notprvNot,notprvAnd,notprvOr,notprvIf}{}{\\[1em]}
\iftag{prvIf}{
  \begin{tabular}{|cc||c|} \hline
    $!A$ & $!B$ & $!A \lif !B$ \\
    \hline \hline
    $\True$ & $\True$ & $\True$ \\
    $\True$ & $\False$ & $\False$ \\
    $\False$ & $\True$ & $\True$ \\
    $\False$ & $\False$ & $\True$ \\
    \hline
  \end{tabular}
}{}
\iftag{prvIff}{
  \begin{tabular}{|cc||c|} \hline
    $!A$ & $!B$ & $!A \liff !B$ \\
    \hline \hline
    $\True$ & $\True$ & $\True$ \\
    $\True$ & $\False$ & $\False$ \\
    $\False$ & $\True$ & $\False$ \\
    $\False$ & $\False$ & $\True$ \\
    \hline
  \end{tabular}
}{}
\end{center}
\end{explain}

\begin{prob}
$\Lang{L_0}$ وچ اک سہ رُکنی رابطہ~$\diamondsuit$ شامل کرن اُتے غور
کرو، جیہدی قدر اگلے طریقے نال دتی گئی اے:
\begin{gather*}
  \pValue{v}(\diamondsuit( !A , !B , !C ) ) = \begin{cases}
    \pValue{v}( !B ) &
    \text{جے $\pValue{v}(!A) = \True$؛}\\
    \pValue{v}( !C ) &
    \text{جے $\pValue{v}(!A) = \False $}.
  \end{cases}
\end{gather*}
ایس رابطے دا صداقتی جدول لکھو۔
\end{prob}

\begin{thm}[مقامی تعیین]
\ollabel{thm:LocalDetermination} منّو کہ $\pAssign{v_1}$ تے
$\pAssign{v_2}$ ایہو جہیاں !!{valuation}s نیں جیہڑیاں~$!A$ وچ آون
والے !!{propositional
variable}s اُتے متفق نیں، یعنی
$\pAssign{v_1}(\Obj p_n) = \pAssign{v_2}(\Obj p_n)$ جدوں وی
$\Obj p_n$، !!{formula}~$!A$ وچ آندا اے۔ تاں $\pValue{v_1}$ تے
$\pValue{v_2}$،~$!A$ اُتے وی متفق نیں، یعنی
$\pValue{v_1}(!A) = \pValue{v_2}(!A)$۔
\end{thm}

\begin{proof}
$!A$ اُتے استقرا نال۔
\end{proof}

\begin{defn}[پورا اترنا]
\ollabel{defn:satisfaction} اسیں
  !!a{formula}~$!A$ دے !!a{valuation}~$\pAssign{v}$ نال
  \emph{پورا اترن} دے تصور، $\pSat{v}{!A}$، دی استقرائی تعریف ایویں
  کر سکدے آں۔ (اسیں $\pSat/{v}{!A}$ نوں ``$\pSat{v}{!A}$ نہیں''
  دے معنی وچ لکھدے آں۔)
\begin{enumerate}
\tagitem{prvFalse}{%
  \indcase{!A}{\lfalse}{$\pSat/{v}{\indfrm}$۔}}{}

\tagitem{prvTrue}{%
  \indcase{!A}{\ltrue}{$\pSat{v}{\indfrm}$۔}}{}

\item \indcase{!A}{\Obj p_i}{$\pSat{v}{\indfrm}$
  اُدوں تے صرف اُدوں جدوں $\pAssign{v}(\Obj p_i) = \True$۔}

\tagitem{prvNot}{%
  \indcase{!A}{\lnot !B}{$\pSat{v}{\indfrm}$ اُدوں تے صرف اُدوں جدوں
    $\pSat/{v}{!B}$۔}}{}

\tagitem{prvAnd}{%
  \indcase{!A}{(!B \land !C)}{$\pSat{v}{\indfrm}$ اُدوں تے صرف اُدوں جدوں
    $\pSat{v}{!B}$ تے $\pSat{v}{!C}$۔}}{}

\tagitem{prvOr}{%
  \indcase{!A}{(!B \lor !C)}{$\pSat{v}{\indfrm}$ اُدوں تے صرف اُدوں جدوں
    $\pSat{v}{!B}$ یا $\pSat{v}{!C}$ (یا دونیں)۔}}{}

\tagitem{prvIf}{%
  \indcase{!A}{(!B \lif !C)}{$\pSat{v}{\indfrm}$ اُدوں تے صرف اُدوں جدوں
    $\pSat/{v}{!B}$ یا $\pSat{v}{!C}$ (یا دونیں)۔}}{}

\tagitem{prvIff}{%
  \indcase{!A}{(!B \liff !C)}{$\pSat{v}{\indfrm}$ اُدوں تے صرف اُدوں جدوں
    یا تاں $\pSat{v}{!B}$ تے $\pSat{v}{!C}$ دونیں، یا نہ
    $\pSat{v}{!B}$ تے نہ $\pSat{v}{!C}$۔}}{}
\end{enumerate}
جے $\Gamma$، !!{formula}s دا سیٹ اے، تاں $\pSat{v}{\Gamma}$ اُدوں تے
صرف اُدوں جدوں $\pSat{v}{!A}$ ہر~$!A \in \Gamma$ لئی۔
\end{defn}

\begin{prop}\ollabel{prop:sat-value}
  $\pSat{v}{!A}$ اُدوں تے صرف اُدوں جدوں $\pValue{v}(!A) = \True$۔
\end{prop}

\begin{proof}
  $!A$ اُتے استقرا نال۔
\end{proof}

\begin{prob}
  \olref[pl][syn][val]{prop:sat-value} ثابت کرو۔
\end{prob}

\end{document}
